\section{Recall Calibration Results}
\label{sec:calibration}

\subsection{The questions-per-prompt effect}
\label{sec:q-effect}

The most surprising finding: asking more questions at once dramatically
improves recall.

\begin{figure}[h]
\centering
\includegraphics[width=0.85\linewidth]{fig1_recall_vs_density.png}
\caption{Recall as a function of fact density (facts/kTok) for three values of $Q$.
The secondary axis maps density values to the \code{d**} notation used throughout the paper.
Grep upper bound (dashed) shows near-perfect keyword presence regardless of density.}
\label{fig:recall-density}
\end{figure}

At $\delta=0.42$ (80 facts in 190 kTok, R4 mode), recall ranges from 67.5\% ($Q=1$) to 78.8\% ($Q=10$) --- an
11pp gap on the same context with the same facts. At higher densities ($\delta=0.63$),
the gap reaches 20pp. Critically, this effect \textbf{reverses} under severe
compaction (\S\ref{sec:q-inversion}): at C4 (98\% compacted), $Q=1$ outperforms $Q=5$ by 9pp, as
focused single-question attention outperforms multi-query probing in degraded contexts.

\textbf{Why this matters for benchmarks}: Any compaction evaluation that uses a
fixed $Q$ is measuring a confound of retrieval ability and multi-query
prompting. Results from different values of $Q$ are not directly comparable.
Our preliminary experiments used $Q=10$ exclusively --- this produced
inflated baselines that masked the true difficulty of recall.

\textbf{Mechanism hypothesis}: Multiple questions in a single prompt create
implicit cross-references that help the model locate relevant information.
A question about ``Rachel's new city'' might prime attention for nearby facts
about ``family trip'' or ``moving costs,'' improving recall for co-located facts.

\subsection{Category hierarchy}
\label{sec:category-hierarchy}

Not all facts are equally recallable. We identify a clear hierarchy:

\begin{figure}[h]
\centering
\includegraphics[width=0.85\linewidth]{fig2_recall_by_category.png}
\caption{Recall breakdown by fact category at $\delta=0.42$ (80 facts), for three values of $Q$.
User statements and assistant responses are reliably recalled (85--100\%),
while preferences remain fragile (20--40\%) regardless of $Q$.}
\label{fig:recall-category}
\end{figure}

\begin{table}[h]
\centering
\caption{Recall hierarchy across LongMemEval fact categories ($\delta=0.42$, $Q=5$).}
\label{tab:category-hierarchy}
\begin{tabular}{lcl}
\toprule
Category & Recall range ($Q=5$) & Notes \\
\midrule
single-session-user       & 80--100\% & User's own statements, easiest \\
single-session-assistant  & 80--100\% & Assistant's responses \\
knowledge-update          & 60--65\%  & Updated information, stable \\
single-session-preference & 20--40\%  & User preferences, fragile \\
temporal-reasoning        & max 30\%  & Requires inference, near-impossible \\
multi-session             & max 42\%  & Cross-session facts, very hard \\
\bottomrule
\end{tabular}
\end{table}

The top two categories (user/assistant statements) are near-ceiling. Knowledge
updates plateau at $\sim$60\%. Preferences are fragile --- the model struggles with
``What is my favorite X?'' even when the information is present. Temporal
reasoning and multi-session facts are effectively unmeasurable in a single
context.

\textbf{Implication for compaction benchmarks}: Testing on ``easy'' categories
(single-session) inflates recall and masks real degradation. Testing on ``hard''
categories (temporal, multi-session) produces floor effects that make strategies
indistinguishable. The filtered mode (R3/R4) excludes impossible categories
while retaining a mix of easy and hard categories.

\subsection{The grep-LLM gap: present but ignored}
\label{sec:grep-llm-gap}

For every fact, we check whether its keywords appear in the context via grep.
Grep recall at $\delta=0.42$ is 86\%. But LLM recall at $Q=1$ is only 67.5\%.

This 19pp gap --- facts that are \textit{verifiably present} in the context but
\textit{not retrieved} by the model --- is a direct measurement of the ``Lost in the
Middle'' phenomenon \citep{liu2023lostmiddle} in a realistic conversational setting.

The gap is not uniform across the context: the spatial recall density
(\Cref{fig:spatial-recall}, \S\ref{sec:spatial}) shows that facts near the beginning and end of the 190K
context are recalled at higher rates than facts in the middle, consistent
with the U-shaped Lost-in-the-Middle pattern reported in the literature.

\subsection{Factorial analysis: filtering vs truncation}
\label{sec:factorial-analysis}

The 2$\times$2 factorial design reveals two independent effects:

\begin{table}[h]
\centering
\caption{Decomposition of the 2$\times$2 factorial effects on recall.}
\label{tab:factorial}
\begin{tabular}{lcl}
\toprule
Effect      & Magnitude & Mechanism \\
\midrule
Filtering   & +14--15pp  & Excludes impossible questions, consistent across truncation \\
Truncation  & $-13$ to $-22$pp  & Removes information from evidence messages \\
Composition & +2.5pp     & Removing hard-category evidence barely helps easy-category recall \\
\bottomrule
\end{tabular}
\end{table}

The \textbf{filtering effect} is remarkably consistent (+14pp and +15pp across
truncation conditions). This means it's driven entirely by excluding
impossible questions, not by freeing context space.

The \textbf{composition effect} is near-zero (+2.5pp mean): whether hard-category
evidence occupies context space alongside easy-category facts makes almost
no difference to easy-category recall. This validates using unfiltered
evidence as realistic padding in controlled experiments.
